\chapter{From Hadronic Regge Trajectories to Light-Cone Strings}

\lectureoverview{String theory did not begin as a theory of quantum gravity.
It began as an attempt to understand the spectrum and scattering of hadrons.
Straight Regge trajectories suggested an extended object; duality in
pion--pion scattering suggested that two apparently distinct processes were
two slicings of one surface. We shall follow those clues into the world-sheet
picture, reinterpret the hadronic string as a QCD flux tube, and then use the
infinite-momentum frame to derive the transverse mechanics of an open string.
The central result is that its effective internal Hamiltonian measures
mass squared. A Hooke-law energy proportional to \(L^2\) in that description
therefore becomes a rest energy proportional to \(L\), the defining behavior
of a string with constant tension.}

\section{The hadronic origin of string theory}

Around 1968--1970, the problem that led to strings concerned protons,
neutrons, and especially mesons. The quark model was already a serious idea,
but the quark content of hadrons had not yet been experimentally established
in the modern sense. Gluons were even less prominent in ordinary discussion.
Nambu had proposed related gauge-field ideas, but the settled language of
quantum chromodynamics did not yet exist. The raw material was instead an
unexpectedly rich spectrum of strongly interacting particles.

It is useful to ignore the proton--neutron distinction temporarily and call
both states the nucleon. Beyond the nucleon were heavier states with larger
spin, followed by still heavier states with still larger spin. Chew and
Frautschi organized such families on a plot with angular momentum on the
vertical axis and mass squared on the horizontal axis. In the historical
units of the period the proton mass was about one BeV; the same unit is now
called a GeV.

Measure angular momentum in units of \(\hbar\). For a baryon family, the
dimensionless spin values are half-integral,

\begin{equation}
  J=\frac12,\frac32,\frac52,\ldots,
\end{equation}

while meson families have integral spin. The remarkable empirical fact was
that either family lay approximately on a straight line when plotted against
\(M^2\). In modern notation one writes

\begin{equation}
  J\simeq \alpha_0+\alpha' M^2,
  \label{eq:l01-regge-trajectory}
\end{equation}

where \(\alpha_0\) is an intercept and \(\alpha'\) is the Regge slope.
\editorialnote{The lecture emphasizes the observed linear relation and its
nearly universal slope. The intercept notation is supplied only to state that
relation compactly; the trajectories are empirical and approximate.}

The original straight line was not compelling because it passed through only
two known points. Any two points determine a line. The surprise came later:
newly discovered states continued to land close to the same line. Meson
trajectories showed the same organization. The pion begins near \(J=0\) with
a small mass on the GeV scale; the rho begins at \(J=1\); further excitations
continue upward. Five, six, or seven states could appear in a family. Even
more strikingly, different baryon and meson trajectories were nearly
parallel. Thus an increase by one quantum of angular momentum corresponded
to approximately the same increase in mass squared:

\begin{equation}
  \Delta J=1
  \quad\Longrightarrow\quad
  \Delta M^2\simeq \frac{1}{\alpha'}.
  \label{eq:l01-regge-spacing}
\end{equation}

\begin{classroomqa}
\audiencequestion{What does it mean for different Regge trajectories to be
parallel?\lecturetimestamp{00:09:32}}
\lecturerresponse{Choose one family of particles and plot all of its excited
states; they lie along one line. Repeat this for another meson or baryon
family. The lines have approximately the same slope even though their
intercepts differ. The same increment of angular momentum therefore costs
approximately the same increment of \(M^2\). This common slope was called the
universal Regge slope.}
\end{classroomqa}

The comparison with ordinary extended objects shows why the pattern mattered.
A basketball acquires rotational energy as it is spun, but its trajectory is
curved rather than linear, and at sufficiently high angular velocity the ball
is torn apart. An atom also has rotational excitations, but eventually it is
ionized; its trajectory ends. Hadrons displayed a much simpler sequence and,
within the available data, no sign that the sequence was giving out.

This spectrum said that a hadron was not behaving like a point. A point has
no spatial structure to spin or deform. An electron has no observed tower of
higher-spin internal states and is correspondingly pointlike within present
experimental bounds. A hadron, by contrast, can be spun up. Its spectrum
suggests compositeness, and the relation between spin and mass hints that the
object deforms as its angular momentum increases.

\section{Duality in pion--pion scattering}

Spectroscopy supplied one clue. Pion--pion scattering supplied another. Two
pions can combine into a rho-like resonance and later emerge again as two
pions. In a conventional Feynman diagram this is an annihilation and
re-creation process, customarily called an \(s\)-channel process. Turning the
same elementary diagram sideways gives a crossed, or \(t\)-channel, process
in which a rho-like particle is exchanged between the pions.

The rho is not alone. It belongs to a Regge family, so consistency appears to
require the rho and all of its excited partners in either channel. Schematically,

\begin{equation}
  \mathcal A(s,t)\simeq
  \sum_n \mathcal A^{(s)}_{\rho_n}(s,t),
  \qquad
  \mathcal A(s,t)\simeq
  \sum_n \mathcal A^{(t)}_{\rho_n}(s,t).
  \label{eq:l01-channel-sums}
\end{equation}

These expressions are not two approximations that should simply be added.
Numerical studies found that summing the tower of resonances in the
annihilation channel already gave approximately the contribution expected
from the exchange channel. Conversely, the exchanged tower could reproduce
the resonance description. Adding both complete towers therefore appeared to
count the same physics twice. Ordinary perturbative bookkeeping says that if
two diagrams exist, both must be included. The hadronic amplitudes instead
suggested that the two channel descriptions were dual.

\begin{figure}[H]
  \centering
  \includegraphics[width=0.82\textwidth]{lecture_01_figure_01.png}
  \caption{The board joins the Chew--Frautschi plot to the annihilation and
  crossed-channel diagrams, making the spectral and scattering clues part of
  one problem.}
  \lectureframe{00:14:20}
  \label{fig:l01-regge-duality-board}
\end{figure}

The same relations can be separated into a clean schematic. The first panel
records the linear spectrum; the other two represent the alternative channel
expansions of the same amplitude.

\begin{figure}[H]
  \centering
  \resizebox{\textwidth}{!}{%
  \begin{tikzpicture}[>=Latex,thick]
    \begin{scope}
      \draw[->] (0,0) -- (4.2,0) node[right] {$M^2$};
      \draw[->] (0,0) -- (0,3.4) node[above] {$J$};
      \draw[blue!65!black] (0.45,0.35) -- (3.55,2.95);
      \draw[orange!80!black] (0.75,0.20) -- (3.85,2.80);
      \foreach \x/\y in {0.65/0.52,1.35/1.10,2.05/1.69,2.75/2.27,3.45/2.86}
        \fill[blue!65!black] (\x,\y) circle (0.07);
      \foreach \x/\y in {0.95/0.37,1.65/0.95,2.35/1.54,3.05/2.12,3.75/2.71}
        \fill[orange!80!black] (\x,\y) circle (0.07);
      \node at (2.1,-0.55) {parallel Regge families};
    \end{scope}
    \begin{scope}[xshift=5.7cm]
      \draw[->] (-1.5,1.2) -- (-0.3,0.25);
      \draw[->] (-1.5,-1.2) -- (-0.3,-0.25);
      \draw[decorate,decoration={snake,amplitude=1.5pt,segment length=7pt}]
        (0,0.25) -- (0,-0.25);
      \draw[->] (0.3,0.25) -- (1.5,1.2);
      \draw[->] (0.3,-0.25) -- (1.5,-1.2);
      \node at (0,-1.75) {$s$-channel tower};
    \end{scope}
    \begin{scope}[xshift=10.5cm]
      \draw[->] (-1.5,1.2) -- (1.5,1.2);
      \draw[->] (-1.5,-1.2) -- (1.5,-1.2);
      \draw[decorate,decoration={snake,amplitude=1.5pt,segment length=7pt}]
        (0,0.95) -- (0,-0.95);
      \node at (0,-1.75) {$t$-channel tower};
    \end{scope}
  \end{tikzpicture}}
  \caption{Parallel Regge families and the two channel expansions involved in
  hadronic duality.}
  \label{fig:l01-regge-duality-redraw}
\end{figure}

\section{One world sheet, two slicings}

The natural response was to seek one geometrical history that could be read
in either way. Represent a pion by a quark--antiquark pair and let time run
upward. Instead of drawing only one-dimensional particle lines, fill in the
region bounded by the quark trajectories. A horizontal slice through this
two-dimensional history may show two incoming strings joining and later
separating. A different slicing displays a string exchanged from one side to
the other. The two Feynman-channel pictures are no longer independent
mechanisms; they are different decompositions of one surface.

\begin{figure}[H]
  \centering
  \includegraphics[width=0.78\textwidth]{lecture_01_figure_02.png}
  \caption{The quark-line topology on the board. Different cuts through the
  same two-dimensional history produce the resonance and exchange readings.}
  \lectureframe{00:17:50}
  \label{fig:l01-worldsheet-topology-board}
\end{figure}

The region between a quark and an antiquark supplies one further idea. At a
fixed time, the slice contains two endpoints and something filling the whole
interval between them. That object is one-dimensional: a string. During an
interaction the string sweeps out the two-dimensional surface.

\begin{definition}
The spacetime history of a point particle is a \emph{world line}. The
spacetime history of a string is a \emph{world sheet}. A string at one instant
is a spatial slice of its world sheet.
\end{definition}

\begin{figure}[H]
  \centering
  \includegraphics[width=0.78\textwidth]{lecture_01_figure_03.png}
  \caption{The board surface filled in between the endpoint trajectories. A
  time slice intersects it in a one-dimensional string.}
  \lectureframe{00:19:15}
  \label{fig:l01-worldsheet-filled-board}
\end{figure}

\begin{figure}[H]
  \centering
  \begin{tikzpicture}[>=Latex,scale=0.95]
    \fill[blue!10]
      (-2.6,-2.3) .. controls (-1.0,-1.0) and (-1.0,-0.4) .. (-0.2,0)
      .. controls (-1.0,0.5) and (-1.0,1.0) .. (-2.6,2.3)
      -- (2.6,2.3)
      .. controls (1.0,1.0) and (1.0,0.5) .. (0.2,0)
      .. controls (1.0,-0.4) and (1.0,-1.0) .. (2.6,-2.3) -- cycle;
    \draw[blue!65!black]
      (-2.6,-2.3) .. controls (-1.0,-1.0) and (-1.0,-0.4) .. (-0.2,0)
      .. controls (-1.0,0.5) and (-1.0,1.0) .. (-2.6,2.3);
    \draw[blue!65!black]
      (2.6,-2.3) .. controls (1.0,-1.0) and (1.0,-0.4) .. (0.2,0)
      .. controls (1.0,0.5) and (1.0,1.0) .. (2.6,2.3);
    \draw[red!75!black,dashed] (-2.9,-1.15) -- (2.9,-1.15)
      node[right] {time slice};
    \draw[orange!80!black,dashed] (-1.55,-2.5) -- (1.55,2.5)
      node[above] {crossed slicing};
    \draw[->] (3.45,-1.7) -- (3.45,1.8) node[above] {time};
    \node[blue!65!black] at (0,-2.7) {one world sheet};
  \end{tikzpicture}
  \caption{A clean world-sheet schematic. Distinct slicings organize the same
  surface into different intermediate-channel descriptions.}
  \label{fig:l01-worldsheet-redraw}
\end{figure}

This geometric reconstruction was exploratory, not a deduction of the full
theory from first principles. Different people reached the early string
picture in different ways. Its strength was that the same extended object
could address both observations: a string can rotate and deform, and its
world sheet can encode the dual scattering channels.

\begin{classroomqa}
\audiencequestion{Why introduce a string-like bridge rather than an ordinary
particle exchanged between the quarks?\lecturetimestamp{00:21:20}}
\lecturerresponse{At one instant there is not merely a localized exchanged
particle. There is a quark, an antiquark, and structure at all points between
them. That structure might be fundamentally continuous or might be a dense
collection of constituents approximating a continuum; the first picture did
not decide. What mattered was that the whole interval was occupied and that
the observed trajectories showed no tendency to terminate like an ordinary
shoelace with weights on its ends.}
\end{classroomqa}

\begin{classroomqa}
\audiencequestion{Can the Standard Model describe pion scattering?
\lecturetimestamp{00:22:45}}
\lecturerresponse{Yes, although usefulness depends on the energy regime. At
low energy the Standard Model, through QCD and its low-energy descriptions,
accounts for a great deal of pion scattering. At higher energies, where whole
towers of hadronic exchanges matter, direct calculations become more
difficult and a string-like organization can be more useful.}
\end{classroomqa}

\begin{classroomqa}
\audiencequestion{Did the early experiments establish that a Regge trajectory
never ends?\lecturetimestamp{00:23:37}}
\lecturerresponse{No. They established only that the observed sequence did
not yet show an endpoint. The inference of indefinite stretchability was a
suggestive extrapolation, not a measured infinity.}
\end{classroomqa}

\section{Flux tubes and the change of scale}

Modern QCD gives the old picture a field-theoretic interpretation. A quark
and an antiquark act as opposite sources of the gluon field. In ordinary
electrodynamics, electric flux spreads sideways as opposite charges are
separated. The field weakens with distance, so the region between the charges
does not retain a fixed transverse profile. QCD is nonlinear. Its dynamics
can confine the relevant chromoelectric flux to a narrow tube joining the
quark and antiquark.

\begin{figure}[H]
  \centering
  \includegraphics[width=0.78\textwidth]{lecture_01_figure_04.png}
  \caption{Quark and antiquark field lines drawn as a narrow flux tube rather
  than a freely spreading Coulomb field.}
  \lectureframe{00:25:15}
  \label{fig:l01-qcd-flux-board}
\end{figure}

\begin{figure}[H]
  \centering
  \begin{tikzpicture}[>=Latex]
    \fill[orange!14,rounded corners=8pt] (-3.0,-0.55) rectangle (3.0,0.55);
    \draw[orange!80!black,thick,rounded corners=8pt]
      (-3.0,-0.55) rectangle (3.0,0.55);
    \fill[blue!70!black] (-3.35,0) circle (0.17);
    \fill[red!75!black] (3.35,0) circle (0.17);
    \node[below] at (-3.35,-0.22) {$q$};
    \node[below] at (3.35,-0.22) {$\bar q$};
    \foreach \y in {-0.30,0,0.30}
      \draw[->,orange!85!black] (-3.05,\y) -- (3.05,\y);
    \draw[<->] (-2.9,-1.05) -- (2.9,-1.05) node[midway,below] {$L$};
    \node at (0,0.92) {approximately fixed transverse profile};
  \end{tikzpicture}
  \caption{A confined flux tube lengthens without spreading. Fixed energy per
  unit length gives \(E=TL\).}
  \label{fig:l01-qcd-flux-redraw}
\end{figure}

The particle language and field language say the same thing differently. In
the field description, the tube retains roughly the same field strength and
cross-sectional area as its endpoints separate. In a constituent picture,
the stretching energy creates more gluonic structure between the endpoints.
An ordinary rubber band has a fixed number of molecules, so its molecules
separate until the material breaks. Imagine instead that a new molecule were
inserted whenever a gap opened. The density along the band could remain fixed
while the band became longer. That is the intended analogy for the gluon
field: the tube gains length rather than thinning toward rupture.

\begin{classroomqa}
\audiencequestion{Would a string that made up an electron also be composed of
gluons?\lecturetimestamp{00:28:08}}
\lecturerresponse{No such conclusion follows. The hadronic string is a QCD
object on the scale of a proton. A hypothetical fundamental string relevant
to quantum gravity would live near the vastly smaller Planck scale, and there
is no evidence here that an electron is a hadronic gluon string. The striking
fact is that related mathematics arose at both scales, not that the physical
objects are identical.}
\end{classroomqa}

Attempts to turn the hadronic picture into a complete interacting string
theory were promising but did not quite produce the desired theory of
hadrons. The mathematical model repeatedly contained a massless particle of
spin two. In a hadronic theory that state looked like an unwanted prediction;
as a theory of gravity it had exactly the quantum numbers of the graviton.
Joel Scherk and John Schwarz made the decisive reinterpretation: perhaps the
theory was telling us about quantum gravity rather than hadrons. The hadronic
and fundamental applications are physically different, even though the early
hadron problem supplied the route to the mathematics.

\begin{classroomqa}
\audiencequestion{How do the flux lines go from a positive endpoint to a
negative endpoint?\lecturetimestamp{00:31:46}}
\lecturerresponse{They leave one source and enter the other, just as magnetic
field lines run from the north pole to the south pole of a bar magnet. The
opposite endpoint labels express the orientation of the flux; a bar magnet
likewise has two opposite poles rather than two north poles.}
\end{classroomqa}

\begin{classroomqa}
\audiencequestion{Is the string continuous or discrete?
\lecturetimestamp{00:32:45}}
\lecturerresponse{Quantum theory makes that distinction subtle. The
electromagnetic field is represented as a continuous field and also has
discrete quanta. In the same way, a string may be treated continuously while
its excitations are quantized. Both descriptions can be valid, depending on
the question being asked.}
\end{classroomqa}

\begin{classroomqa}
\audiencequestion{If the angle of the spacetime slice through a world sheet
is changed, is the intersection still a string?\lecturetimestamp{00:34:20}}
\lecturerresponse{Yes. Changing the simultaneity slice is a Lorentz
transformation. The new slice intersects the surface in a string seen in a
different state of motion. Even a string whose center of mass is at rest is
not internally motionless; it may vibrate. ``At rest'' refers only to its
center of mass.}
\end{classroomqa}

\section{Two expansions of relativistic energy}

A direct relativistic description of an extended, vibrating object is
awkward. The useful strategy is to find a frame in which its relevant motion
has the form of ordinary nonrelativistic mechanics. Begin with a single
particle. Nonrelativistically one writes

\begin{equation}
  E_{\mathrm{nr}}=\frac{\bm p^2}{2m}+B,
  \label{eq:l01-nonrel-energy}
\end{equation}

where \(B\) is an internal or binding energy independent of the overall
momentum. For several noninteracting particles, both the kinetic terms and
the internal constants are added. Since only energy differences govern the
motion, a state-independent additive constant can often be suppressed.

The exact relativistic energy is

\begin{equation}
  E=\sqrt{\bm p^2c^2+m^2c^4}.
  \label{eq:l01-rel-energy-c}
\end{equation}

\begin{figure}[H]
  \centering
  \includegraphics[width=0.76\textwidth]{lecture_01_figure_05.png}
  \caption{The nonrelativistic kinetic energy and the exact relativistic
  square root, with the powers of \(c\) corrected during the discussion.}
  \lectureframe{00:38:45}
  \label{fig:l01-relativistic-energy-board}
\end{figure}

\begin{classroomqa}
\audiencequestion{Should the relativistic formula contain \(p^2c^2\) and
\(m^2c^4\), rather than the powers initially written on the board?
\lecturetimestamp{00:39:12}}
\lecturerresponse{Yes. Dimensional consistency fixes the expression to
\(E^2=p^2c^2+m^2c^4\). After recovering the nonrelativistic limit we shall set
\(c=1\), but keeping the powers here makes the expansion transparent.}
\end{classroomqa}

When \(p^2c^2\ll m^2c^4\), factor out \(m^2c^4\) and use
\(\sqrt{1+\epsilon}\simeq1+\epsilon/2\):

\begin{align}
  E
  &=mc^2\sqrt{1+\frac{p^2}{m^2c^2}}
  \notag\\
  &\simeq mc^2+\frac{p^2}{2m}.
  \label{eq:l01-low-momentum-expansion}
\end{align}

\begin{figure}[H]
  \centering
  \includegraphics[width=0.76\textwidth]{lecture_01_figure_06.png}
  \caption{The small-momentum expansion recovering rest energy plus the
  familiar nonrelativistic kinetic term.}
  \lectureframe{00:41:30}
  \label{fig:l01-low-momentum-board}
\end{figure}

This approximation requires more than a slowly moving center of mass. A box
may drift slowly while particles inside it move near the speed of light. In
that case the internal dynamics are still relativistic. The ordinary
nonrelativistic limit requires all relevant constituent velocities to be
small in the chosen frame.

There is another useful limit. Boost the entire system along the \(z\)-axis
until every constituent has a large positive longitudinal momentum \(P_{zi}\).
Any constituent initially moving backward is carried forward by a
sufficiently large boost. The transverse momenta and rest masses remain fixed:

\begin{equation}
  P_{zi}\gg |p_{xi}|,|p_{yi}|,m_i,
  \qquad
  p_{\perp i}^2=p_{xi}^2+p_{yi}^2,
\end{equation}

where from this point onward \(c=1\). Expanding for large \(P_{zi}\) gives

\begin{align}
  E_i
  &=\sqrt{P_{zi}^2+p_{\perp i}^2+m_i^2}
  \notag\\
  &=P_{zi}\sqrt{1+\frac{p_{\perp i}^2+m_i^2}{P_{zi}^2}}
  \notag\\
  &\simeq P_{zi}
  +\frac{p_{\perp i}^2+m_i^2}{2P_{zi}}.
  \label{eq:l01-large-pz-one-particle}
\end{align}

\begin{figure}[H]
  \centering
  \includegraphics[width=0.76\textwidth]{lecture_01_figure_07.png}
  \caption{The large-\(P_z\) square-root expansion, with transverse momentum
  and mass retained as the small quantities.}
  \lectureframe{00:46:45}
  \label{fig:l01-large-pz-board}
\end{figure}

For many constituents,

\begin{equation}
  E\simeq
  \sum_i P_{zi}
  +\sum_i\frac{p_{\perp i}^2}{2P_{zi}}
  +\sum_i\frac{m_i^2}{2P_{zi}}.
  \label{eq:l01-large-pz-sum}
\end{equation}

The first term is the conserved total longitudinal momentum. It is enormous,
but it is also a fixed additive contribution. Keeping or subtracting it
changes no energy difference, just as adding any other conserved number to a
Hamiltonian would not change the equations of motion. The dynamical part is
therefore

\begin{equation}
  \begin{aligned}
    H_{\perp}=E-P_z
    &\simeq
    \sum_i\left(
      \frac{p_{\perp i}^2}{2P_{zi}}
      +\frac{m_i^2}{2P_{zi}}
    \right),\\
    P_z&=\sum_iP_{zi}.
  \end{aligned}
  \label{eq:l01-transverse-hamiltonian}
\end{equation}

\begin{figure}[H]
  \centering
  \includegraphics[width=0.78\textwidth]{lecture_01_figure_08.png}
  \caption{The many-particle decomposition into conserved longitudinal
  momentum, transverse kinetic energy, and mass-dependent internal terms.}
  \lectureframe{00:49:16}
  \label{fig:l01-energy-sum-board}
\end{figure}

\section{The infinite-momentum frame}

The preceding limit was historically called the infinite-momentum frame and
is now usually formulated in light-cone coordinates. Its meaning can be read
directly from the Hamiltonian. In quantum mechanics,

\begin{equation}
  i\hbar\frac{\partial}{\partial t}\ket{\psi}=H\ket{\psi}.
  \label{eq:l01-time-evolution}
\end{equation}

All nontrivial terms in \(H_{\perp}\) scale as \(1/P_z\). As the boost grows,
energy splittings shrink and changes occur more slowly. This is time
dilation. A radioactive system appears to decay more slowly when highly
boosted; one can compensate by observing it for a proportionally longer time.
The slowing does not remove the dynamics. It rescales their clock.

\begin{figure}[H]
  \centering
  \includegraphics[width=0.78\textwidth]{lecture_01_figure_09.png}
  \caption{The board identifies the reduced energy with the Hamiltonian that
  generates the slowed time evolution of the boosted system.}
  \lectureframe{00:52:57}
  \label{fig:l01-hamiltonian-board}
\end{figure}

The transverse term has precisely the nonrelativistic form \(p^2/(2m)\), but
the role of mass is played by \(P_z\). This is physically sensible: a larger
longitudinal momentum makes a particle harder to deflect sideways, so it acts
as transverse inertia. The term \(m_i^2/(2P_{zi})\) does not depend on the
two-dimensional state of motion and therefore plays the role of an internal
or binding energy. Relativistic motion in three spatial dimensions has become
nonrelativistic-looking motion in the two dimensions perpendicular to the
boost.

This is the permission needed to use an elementary masses-and-springs model
for a relativistic string. The model is not pretending that the string is
slow in its rest frame. It uses the frame in which the transverse dynamics
have the required nonrelativistic structure.

\section{An open string made from beads and springs}

Take an open string, meaning a string with two endpoints. Something may later
be attached to those endpoints, but the present concern is the string itself.
Approximate it by beads at transverse positions
\(\bm X_i(t)=(X_i(t),Y_i(t))\), joined by identical springs. A discrete
Hamiltonian is

\begin{equation}
  H_N=\sum_{i=1}^{N}
  \left[
    \frac{\mu_N}{2}\dot{\bm X}_i^{\,2}
    +\frac{k_N}{2}
      \left(\bm X_{i+1}-\bm X_i\right)^2
  \right].
  \label{eq:l01-discrete-string}
\end{equation}

The kinetic term alone would let the beads fly apart. The nearest-neighbor
Hooke term is what makes them a string.

\begin{figure}[H]
  \centering
  \includegraphics[width=0.75\textwidth]{lecture_01_figure_10.png}
  \caption{The first bead-and-spring sketch for an open string with two
  transverse coordinates.}
  \lectureframe{01:00:30}
  \label{fig:l01-bead-chain-board}
\end{figure}

\begin{figure}[H]
  \centering
  \begin{tikzpicture}[>=Latex]
    \foreach \x in {0,...,8} {
      \fill[blue!65!black] ({0.85*\x},0) circle (0.09);
    }
    \foreach \x in {0,...,7} {
      \draw[decorate,decoration={zigzag,segment length=4pt,amplitude=2pt}]
        ({0.85*\x+0.10},0) -- ({0.85*\x+0.75},0);
    }
    \draw[->] (3.4,0.15) -- (3.4,1.25) node[above] {$\bm X_i(t)$};
    \draw[<->] (0,-0.65) -- (6.8,-0.65)
      node[midway,below] {$0\leq\sigma\leq\pi$};
  \end{tikzpicture}
  \caption{Discrete regulator for the open string. The bead mass and spring
  constant scale with the spacing before the continuum limit is taken.}
  \label{fig:l01-bead-chain-redraw}
\end{figure}

\begin{figure}[H]
  \centering
  \includegraphics[width=0.78\textwidth]{lecture_01_figure_11.png}
  \caption{The discrete kinetic and Hooke terms written together before the
  continuum limit.}
  \lectureframe{01:01:30}
  \label{fig:l01-discrete-energy-board}
\end{figure}

As the number of beads increases, each bead mass \(\mu_N\) must tend to zero
so that the whole string retains finite inertia. At the same time the spring
constant \(k_N\) must grow. A long strip of rubber is easier to stretch by a
fixed amount than a tiny segment of the same material; neighboring points
that are brought closer together must therefore be connected more stiffly if
the continuum material is to remain the same.

Introduce a coordinate \(\sigma\) along the string. Its numerical range is a
convention. Choosing

\begin{equation}
  0\leq\sigma\leq\pi
\end{equation}

is convenient because a closed string can later be parameterized from
\(0\) to \(2\pi\). The sum becomes an integral, and the nearest-neighbor
difference becomes a derivative:

\begin{equation}
  \sum_i\longrightarrow\int_0^\pi d\sigma,
  \qquad
  \bm X_{i+1}-\bm X_i
  \longrightarrow
  \partial_\sigma\bm X.
  \label{eq:l01-continuum-replacements}
\end{equation}

After absorbing the spacing-dependent constants into the normalization of
\(t\), \(\sigma\), and \(\bm X\), the continuum Hamiltonian has the form

\begin{equation}
  H_{\mathrm{string}}=
  \frac12\int_0^\pi d\sigma
  \left[
    \left(\partial_t\bm X\right)^2
    +\left(\partial_\sigma\bm X\right)^2
  \right].
  \label{eq:l01-continuum-hamiltonian}
\end{equation}

\begin{figure}[H]
  \centering
  \includegraphics[width=0.78\textwidth]{lecture_01_figure_12.png}
  \caption{The board passage from the bead sum to
  \(\int_0^\pi d\sigma\), with the nearest-neighbor difference replaced by
  \(\partial_\sigma\bm X\).}
  \lectureframe{01:05:00}
  \label{fig:l01-continuum-board}
\end{figure}

The first term is kinetic energy in the two transverse directions. The second
measures stretching along the string. Both terms carry a plus sign because
this is the Hamiltonian. In the Lagrangian the potential term would enter
with a minus sign. The precise coefficients depend on conventions and were
not fixed at this stage; the kinetic-plus-stretching structure is what will
matter.

Separate the center-of-mass motion from the relative motion. Translating
every point by the same amount moves the particle as a whole. Oscillations,
bending, and relative stretching are internal motions. A proton, an atom, or
any other composite object can be called a particle even though its
constituents move internally. The string is no different in that respect.

\section{Why the Hamiltonian measures mass squared}

For a slowly moving object, rest energy is proportional to its mass. In the
infinite-momentum description, however, the internal contribution appears as
\(M^2/(2P_z)\). Once the common longitudinal momentum and its associated time
rescaling have been fixed, the internal string Hamiltonian is proportional to
the invariant mass squared of the whole object:

\begin{equation}
  H_{\mathrm{int}}\propto M^2.
  \label{eq:l01-hint-mass-squared}
\end{equation}

\editorialnote{More explicitly, light-cone kinematics gives an invariant
relation of the form \(M^2=2P^+P^- - \bm P_\perp^2\). In the center-of-mass
transverse frame and at fixed \(P^+\), the light-cone Hamiltonian \(P^-\) is
therefore proportional to \(M^2\). The lecture suppresses these normalization
details and keeps the proportionality needed for the argument.}

\begin{figure}[H]
  \centering
  \includegraphics[width=0.76\textwidth]{lecture_01_figure_13.png}
  \caption{The board identifies the internal string energy with the mass
  squared of the complete string state.}
  \lectureframe{01:10:30}
  \label{fig:l01-mass-squared-board}
\end{figure}

A vibrating string is a collection of coupled oscillators. Its complete
quantization is postponed, but the important spectral fact is already
visible. Harmonic-oscillator energies rise in discrete steps. Since those
internal energies contribute to \(M^2\), string excitations naturally produce
discrete steps in mass squared. This is the structure that the Regge plots
announced at the beginning.

\begin{classroomqa}
\audiencequestion{Is \(M\) the mass of the whole string, and would doubling
that mass quadruple the corresponding \(M^2\) contribution?
\lecturetimestamp{01:12:10}}
\lecturerresponse{Yes. \(M\) is the rest mass of the complete string state,
including vibration, stretching, and every other internal contribution. The
effective boosted Hamiltonian is proportional to \(M^2\), so a state with
twice the mass has four times that mass-squared value.}
\end{classroomqa}

\begin{classroomqa}
\audiencequestion{How can this description include a photon, whose rest mass
is zero?\lecturetimestamp{01:12:52}}
\lecturerresponse{A massless state corresponds to \(M=0\). How the positive
oscillator contributions can combine to produce such a state is a
significant question for the quantized theory, not something settled by the
classical model here. At this point \(M\) simply denotes the total rest mass,
and quantum mechanics has not yet been introduced.}
\end{classroomqa}

Now stretch a nonmoving string uniformly to physical length \(L\). The
kinetic term vanishes, and over the parameter interval of length \(\pi\),

\begin{equation}
  \left|\partial_\sigma\bm X\right|\simeq\frac{L}{\pi}.
\end{equation}

Consequently,

\begin{align}
  H_{\mathrm{stretch}}
  &\propto
  \int_0^\pi d\sigma
    \left(\partial_\sigma\bm X\right)^2
  \notag\\
  &\simeq
  \pi\left(\frac{L}{\pi}\right)^2
  \propto L^2.
  \label{eq:l01-hooke-l2}
\end{align}

\begin{figure}[H]
  \centering
  \includegraphics[width=0.76\textwidth]{lecture_01_figure_14.png}
  \caption{The uniformly stretched configuration used to infer
  \(H_{\mathrm{stretch}}\propto L^2\) and hence \(M\propto L\).}
  \lectureframe{01:15:30}
  \label{fig:l01-length-relation-board}
\end{figure}

This is Hooke-law behavior in the effective transverse mechanics. But the
quantity on the left is proportional to \(M^2\), not \(M\). It follows that

\begin{equation}
  M^2\propto L^2,
  \qquad
  M\propto L.
  \label{eq:l01-mass-length}
\end{equation}

In the rest frame the energy is \(E_{\mathrm{rest}}=M\), in units \(c=1\).
Writing the constant of proportionality as the string tension \(T\),

\begin{equation}
  E_{\mathrm{rest}}=TL.
  \label{eq:l01-string-tension}
\end{equation}

The result matches the flux-tube picture. A uniform electric or magnetic
field carries a uniform energy density. If its flux is confined to a tube of
fixed cross section, each additional unit of length adds the same amount of
energy. Equivalently, a constituent picture in which new degrees of freedom
fill each newly opened gap has fixed energy per unit length. String tension is
that energy per unit length.

This conclusion is indirect. It does not claim that an arbitrary rubber band
driven to relativistic speeds will automatically acquire \(E=TL\). The
argument begins with the transverse nonrelativistic description that follows
from the large boost, identifies its internal energy with \(M^2\), and only
then translates the result back to rest-frame energy.

\section{The same object in two frames}

The distinction between \(M\) and \(M^2\) is easy to lose, so it is worth
placing the two expansions side by side. With \(c=1\),

\begin{align}
  E&=\sqrt{p^2+M^2}
  \simeq M+\frac{p^2}{2M},
  &&p\ll M,
  \label{eq:l01-rest-expansion-again}\\
  E&=\sqrt{P_z^2+M^2}
  \simeq P_z+\frac{M^2}{2P_z},
  &&P_z\gg M.
  \label{eq:l01-boost-expansion-again}
\end{align}

\begin{figure}[H]
  \centering
  \includegraphics[width=0.74\textwidth]{lecture_01_figure_15.png}
  \caption{The small-momentum expansion repeated after the break, where
  internal rest energy is read as \(M\).}
  \lectureframe{01:28:00}
  \label{fig:l01-rest-expansion-board-again}
\end{figure}

\begin{figure}[H]
  \centering
  \includegraphics[width=0.74\textwidth]{lecture_01_figure_16.png}
  \caption{The complementary large-momentum expansion, where the residual
  energy is proportional to \(M^2/(2P_z)\).}
  \lectureframe{01:31:00}
  \label{fig:l01-boost-expansion-board-again}
\end{figure}

The first expansion identifies rest energy with mass and ordinary inertia
with \(M\). In the second, the conserved \(P_z\) term is removed, transverse
inertia is \(P_z\), and internal energy is represented by \(M^2/(2P_z)\). The
mass remains proportional to the string length in either frame. What changes
is which function of the mass appears as the residual Hamiltonian.

\begin{classroomqa}
\audiencequestion{A Lorentz boost should not change the physics. Why does the
boosted string appear to obey a different energy law?
\lecturetimestamp{01:25:35}}
\lecturerresponse{It does not obey different physics. In both descriptions
\(M\propto L\) and \(M^2\propto L^2\). In the rest frame the relevant energy
is \(M\); after a large boost the conserved \(P_z\) is separated off and the
remaining Hamiltonian is proportional to \(M^2\). The boost changes the useful
description, not the invariant mass--length relation.}
\end{classroomqa}

\section{A superconducting model of confinement}

A superconductor expels magnetic field from its bulk. If magnetic flux is
trapped through a sufficiently narrow channel before the material becomes
superconducting, the flux cannot spread sideways and remains confined to a
tube. A useful Gedanken experiment is to drill a narrow hole through a
normal-state sample, insert a long bar magnet through it, cool the sample, and
then withdraw the magnet. The construction is schematic rather than a
practical laboratory recipe. Its purpose is to isolate a flux line whose
field strength and energy density are approximately uniform along its length.

\begin{figure}[H]
  \centering
  \includegraphics[width=0.76\textwidth]{lecture_01_figure_17.png}
  \caption{The superconducting flux-line construction and the corresponding
  monopole--antimonopole confinement picture.}
  \lectureframe{01:35:30}
  \label{fig:l01-superconductor-board}
\end{figure}

Suppose, as a further thought experiment, that a magnetic monopole and
antimonopole could be placed in the superconducting medium and separated. The
magnetic flux between them would be squeezed into a tube. Since each unit of
tube length carries the same energy,

\begin{equation}
  V(R)=T_{\mathrm{flux}}R.
  \label{eq:l01-linear-confinement}
\end{equation}

The force \(F=-dV/dR\) has constant magnitude. Pulling the pair apart becomes
increasingly expensive, so the pair is confined. This is the behavior sought
for quarks in a hadron.

\begin{figure}[H]
  \centering
  \begin{tikzpicture}[>=Latex]
    \fill[blue!8,rounded corners=5pt] (-4,-1.35) rectangle (4,1.35);
    \node[blue!55!black] at (0,1.65) {superconducting medium};
    \fill[red!75!black] (-2.8,0) circle (0.17);
    \fill[blue!70!black] (2.8,0) circle (0.17);
    \node[below] at (-2.8,-0.22) {monopole};
    \node[below] at (2.8,-0.22) {antimonopole};
    \fill[orange!18,rounded corners=7pt] (-2.65,-0.42) rectangle (2.65,0.42);
    \draw[orange!85!black,thick,rounded corners=7pt]
      (-2.65,-0.42) rectangle (2.65,0.42);
    \foreach \y in {-0.22,0,0.22}
      \draw[->,orange!85!black] (-2.55,\y) -- (2.55,\y);
    \draw[<->] (-2.8,-1.00) -- (2.8,-1.00) node[midway,below] {$R$};
  \end{tikzpicture}
  \caption{Flux confinement gives a uniform tube and a linear potential
  \(V(R)=T_{\mathrm{flux}}R\).}
  \label{fig:l01-superconductor-redraw}
\end{figure}

\begin{classroomqa}
\audiencequestion{What plays the role of the superconducting medium in the
hadron?\lecturetimestamp{01:36:57}}
\lecturerresponse{The proposed analogy reverses the electric and magnetic
roles. An ordinary superconducting condensate is electrically charged and
confines magnetic flux. In the dual-superconductor picture of QCD, a
condensate of chromomagnetic objects confines the chromoelectric charge of
quarks. Such objects do not appear as isolated particles because they are
represented as condensed in the vacuum.}
\end{classroomqa}

\editorialnote{The dual-superconductor description is an influential
confinement mechanism and is presented here in the historical language of the
lecture. It should be read as a physical picture of the QCD vacuum, not as a
complete derivation of confinement from first principles.}

\begin{classroomqa}
\audiencequestion{Is energy proportional to length a general rule for any
ordinary string made highly relativistic?\lecturetimestamp{01:38:10}}
\lecturerresponse{No. One would first need a consistent relativistic theory
of that material, and its behavior would depend on the details. Linear energy
is special to a configuration whose density stays fixed as it is extended,
or to flux prevented from spreading transversely. Pulling it apart creates
more constituents or more uniform field-filled length rather than merely
increasing the separation of a fixed set of molecules.}
\end{classroomqa}

\section{Transverse degrees of freedom and the size of a particle}

The light-cone construction writes the string directly in terms of its
transverse coordinates. There is no separate, freely specifiable longitudinal
oscillator analogous to each transverse oscillator. This reduction by one
direct spatial dimension is one of the remarkable features of the formalism;
the complete theory nevertheless has to satisfy Lorentz invariance.

\begin{classroomqa}
\audiencequestion{What happens to residual motion along the \(z\)-direction?
\lecturetimestamp{01:40:10}}
\lecturerresponse{In the light-cone description, longitudinal motion is not
introduced as an independent degree of freedom. The physical oscillations
are described transversely, while constraints determine the longitudinal
information. This economy is related in spirit to the idea that gravitating
systems may admit a description with one fewer direct dimension, although
that connection is not derived here.}
\end{classroomqa}

\begin{classroomqa}
\audiencequestion{How many points are actually known on the Regge plots?
\lecturetimestamp{01:41:15}}
\lecturerresponse{The exact maximum depends on the family; a rough number is
about seven. The sequence is not experimentally endless. The relevant fact
is that the known points do not display the kind of visible failure or
breakup endpoint familiar from ordinary macroscopic material.}
\end{classroomqa}

Hadronic phenomenology supplied the historical motivation, but the
mathematical string theory is not required to reproduce hadrons in every
detail. Once the formalism is built, one asks what it actually describes. Its
massless spin-two state and other properties point more naturally toward
gravity.

The apparent contrast between hadrons and electrons can be quantified with
the energy of a rigid rotation,

\begin{equation}
  E_{\mathrm{rot}}=\frac{J^2}{2I},
  \qquad
  I\sim MR^2.
  \label{eq:l01-rotational-energy}
\end{equation}

Angular momentum is quantized. Writing \(J^2\) schematically as
\(\hbar^2 n^2\), the first rotational step is of order

\begin{equation}
  \Delta E_{0\to1}\sim\frac{\hbar^2}{2I}.
  \label{eq:l01-first-rotational-step}
\end{equation}

\begin{figure}[H]
  \centering
  \includegraphics[width=0.76\textwidth]{lecture_01_figure_18.png}
  \caption{The moment-of-inertia estimate for the first rotational
  excitation, \(\Delta E\sim\hbar^2/(2I)\).}
  \lectureframe{01:44:00}
  \label{fig:l01-rotational-board}
\end{figure}

For fixed mass, a smaller radius means a much smaller moment of inertia and a
much larger excitation energy. Hadrons are large on the scale of their
internal dynamics, so a rotational excitation costs roughly a GeV and is
accessible in particle collisions. If electrons possess any spatial
substructure, experimental bounds make its scale far smaller, and rotational
excitations would lie far beyond available energies. A basketball has the
opposite extreme: its moment of inertia is so large that one quantum of
angular momentum costs an imperceptibly small energy. Atomic rotational or
orbital steps lie between these extremes, commonly at the scale of electron
volts.

\begin{classroomqa}
\audiencequestion{Could an apparently elementary electron still have spatial
structure and rotational excitations?\lecturetimestamp{01:42:04}}
\lecturerresponse{It could only hide them at very high energy. The cost of one
unit of angular momentum scales as \(1/I\), and \(I\sim MR^2\). An exceedingly
small radius therefore pushes the first excitation far away on the
mass-squared plot. The absence of observed electron excitations means either
that the electron is pointlike or that any compositeness scale is much smaller
than present experiments can resolve.}
\end{classroomqa}

The opening Regge plot can now be read dynamically. Hadrons reveal accessible
rotational and vibrational states because they are extended. The equal steps
in their mass-squared spectrum are the signature expected when an oscillator
Hamiltonian is identified with \(M^2\). Their approximate linear energy with
length is the signature of a flux tube with constant tension. Spectroscopy,
scattering duality, world-sheet geometry, and light-cone mechanics are four
descriptions of the same organizing idea: an extended relativistic object
whose natural internal variable is mass squared.
